%% ============================================================
%%  Chapter 4 – Experiments and Results
%% ============================================================
\chapter{Experiments and Results}
\label{ch:results}

\begin{infobox}[Chapter Overview]
This chapter describes the experimental setup, dataset statistics,
quantitative evaluation results, ablation studies, and qualitative
case analysis.
\end{infobox}

%% ─────────────────────────────────────────────────────────────
\section{Experimental Setup}
\label{sec:setup}

\subsection{Hardware and Software}
\label{subsec:hw}

All experiments are conducted on a workstation equipped with a single
NVIDIA GPU (recommended: $\geq$16\,GB VRAM), 32\,GB system RAM, and
an NVMe SSD for fast DICOM I/O.
The software stack is:

\begin{itemize}
  \item Python 3.10, PyTorch 2.1~\citep{paszke2019pytorch}
  \item Hugging Face Transformers 4.38~\citep{wolf2020transformers}
  \item pydicom 2.4 + pylibjpeg 2.0 (DICOM I/O)
  \item Gradio 4.x (clinical UI)
  \item bert-score 0.3.13, NLTK 3.8 (evaluation)
\end{itemize}

Reproducibility is ensured by fixing global seeds:
\begin{lstlisting}[caption={Seed setting for reproducibility.}]
import torch, random, numpy as np
torch.manual_seed(42)
np.random.seed(42)
random.seed(42)
\end{lstlisting}

\subsection{Hyperparameters}
\label{subsec:hparams}

Table~\ref{tab:hparams} lists all key hyperparameters.

\begin{table}[htbp]
\centering
\caption{Training hyperparameters.}
\label{tab:hparams}
\begin{tabular}{@{}lll@{}}
\toprule
\textbf{Hyperparameter} & \textbf{Phase 1} & \textbf{Phase 2} \\
\midrule
Epochs              & 10            & 20 \\
Batch size          & 8             & 8 \\
Optimiser           & AdamW         & AdamW \\
Learning rate       & $1\times10^{-4}$ & $5\times10^{-5}$ \\
LR schedule         & Cosine        & Cosine \\
Weight decay        & $1\times10^{-4}$ & $1\times10^{-4}$ \\
Grad. clip norm     & 1.0           & 1.0 \\
AMP                 & \checkmark    & \checkmark \\
LoRA rank           & ---           & 8 \\
LoRA $\alpha$       & ---           & 16 \\
LoRA targets        & ---           & query, key, value \\
NT-Xent temperature & 0.07          & --- \\
Val. every (steps)  & 50            & 50 \\
\bottomrule
\end{tabular}
\end{table}

%% ─────────────────────────────────────────────────────────────
\section{Dataset}
\label{sec:dataset}

\begin{warningbox}[Data Privacy Notice]
All data used in this work was de-identified by the hospital informatics
team using \medxplain's \code{phi\_scrubber} module before being
transferred to the research environment. No PHI is present in any
experimental artefact. Access was granted under an IRB-approved protocol.
\end{warningbox}

\subsection{Cohort Description}
\label{subsec:cohort}

Table~\ref{tab:cohort} summarises the hospital chest radiograph cohort
used in this work. Values marked with ``---'' should be filled in with
actual figures obtained from the hospital informatics team following
IRB approval.

\begin{table}[htbp]
\centering
\caption{Hospital Chest Radiograph Cohort Statistics.}
\label{tab:cohort}
\begin{tabular}{@{}lrrr@{}}
\toprule
\textbf{Split} & \textbf{Studies} & \textbf{QA Pairs} & \textbf{Modalities} \\
\midrule
Train           & ---  & ---  & CR, DX, CT \\
Validation      & ---  & ---  & CR, DX, CT \\
Test (Annotated)& ---  & ---  & CR, DX, CT \\
\midrule
\textbf{Total}  & ---  & ---  & \\
\bottomrule
\end{tabular}
\end{table}

\subsection{Annotation Protocol}
\label{subsec:annotation}

Ground-truth labels were generated by two board-certified radiologists who
independently answered each clinical question. Disagreements were adjudicated
by a senior radiologist. The annotation form (Appendix~\ref{app:form}) was
administered \emph{before} the annotator saw the model output, preventing
anchoring bias.

\subsection{Pathology Distribution}
\label{subsec:distribution}

Table~\ref{tab:findings} shows the per-pathology prevalence in the test set,
as extracted by \code{report\_parser.py} from the annotated radiology reports.

\begin{table}[htbp]
\centering
\caption{Pathology prevalence in the test set.}
\label{tab:findings}
\begin{tabular}{@{}lrr@{}}
\toprule
\textbf{Pathology} & \textbf{N Positive} & \textbf{Prevalence (\%)} \\
\midrule
Pleural Effusion    & --- & --- \\
Cardiomegaly        & --- & --- \\
Pneumonia           & --- & --- \\
Atelectasis         & --- & --- \\
Pneumothorax        & --- & --- \\
Pulmonary Nodule    & --- & --- \\
Pulmonary Oedema    & --- & --- \\
Interstitial Pattern& --- & --- \\
\bottomrule
\end{tabular}
\end{table}

%% ─────────────────────────────────────────────────────────────
\section{Quantitative Results}
\label{sec:quant}

\subsection{Primary Metrics}
\label{subsec:primary}

Table~\ref{tab:main_results} presents the primary evaluation results on
the held-out test set, generated by
\code{evaluation/evaluate\_metrics\_spec.py}.

\begin{table}[htbp]
\centering
\caption{Primary evaluation results vs.\ Cahier des Charges targets.}
\label{tab:main_results}
\begin{tabular}{@{}lrrl@{}}
\toprule
\textbf{Metric}     & \textbf{Score} & \textbf{Target} & \textbf{Status} \\
\midrule
VQA Accuracy        & ---   & $>65\%$  & --- \\
BLEU-4              & ---   & $>0.25$  & --- \\
BERTScore F1        & ---   & $>0.75$  & --- \\
Clinical F1 (Macro) & ---   & $>0.60$  & --- \\
\bottomrule
\end{tabular}
\end{table}

\subsection{Per-Pathology Clinical F1}
\label{subsec:clinical_f1}

Clinical F1 measures whether the model correctly identifies the
presence or absence of each pathology in its free-text answer.
For a pathology $p$, let $\mathrm{TP}_p$, $\mathrm{FP}_p$, $\mathrm{FN}_p$
be the true positives, false positives, and false negatives respectively.
Then:
\begin{equation}
  F1_p = \frac{2\,\mathrm{TP}_p}{2\,\mathrm{TP}_p + \mathrm{FP}_p + \mathrm{FN}_p}
  \label{eq:f1}
\end{equation}
Macro-average Clinical F1 is:
\begin{equation}
  \overline{F1} = \frac{1}{|\mathcal{P}|} \sum_{p \in \mathcal{P}} F1_p
  \label{eq:macro_f1}
\end{equation}
where $\mathcal{P}$ includes only pathologies with at least one positive
reference in the test set.

Table~\ref{tab:per_path} reports per-pathology results.

\begin{table}[htbp]
\centering
\caption{Per-pathology Clinical F1 scores.}
\label{tab:per_path}
\begin{tabular}{@{}lrr@{}}
\toprule
\textbf{Pathology}   & \textbf{F1} & \textbf{Prevalence (\%)} \\
\midrule
Pneumothorax         & --- & --- \\
Pleural Effusion     & --- & --- \\
Pneumonia            & --- & --- \\
Cardiomegaly         & --- & --- \\
Atelectasis          & --- & --- \\
Pulmonary Oedema     & --- & --- \\
Pulmonary Nodule     & --- & --- \\
\midrule
\textbf{Macro Average} & --- & --- \\
\bottomrule
\end{tabular}
\end{table}

\subsection{Radiologist Agreement}
\label{subsec:kappa}

Cohen's Kappa (Equation~\ref{eq:kappa}) between model and radiologist
predictions:

\begin{table}[htbp]
\centering
\caption{Model–radiologist agreement.}
\label{tab:kappa}
\begin{tabular}{@{}lrr@{}}
\toprule
\textbf{Metric}      & \textbf{Value} & \textbf{Interpretation} \\
\midrule
Accuracy             & ---   & Fraction of exact matches \\
Cohen's Kappa        & ---   & Agreement beyond chance \\
N discrepancies      & ---   & Cases flagged for review \\
\bottomrule
\end{tabular}
\end{table}

Kappa values are interpreted as: $\kappa < 0.2$ slight; $0.2$--$0.4$ fair;
$0.4$--$0.6$ moderate; $0.6$--$0.8$ substantial; $>0.8$ almost perfect
agreement~\citep{landis1977kappa}.

%% ─────────────────────────────────────────────────────────────
\section{Ablation Studies}
\label{sec:ablation}

\subsection{Effect of Phase 1 Pre-alignment}
\label{subsec:abl_phase1}

Table~\ref{tab:abl_phase} compares Phase~2 training with and without the
Phase~1 alignment initialisation.

\begin{table}[htbp]
\centering
\caption{Ablation: Impact of Phase 1 vision-language alignment.}
\label{tab:abl_phase}
\begin{tabular}{@{}lccc@{}}
\toprule
\textbf{Configuration} & \textbf{VQA Acc.} & \textbf{BLEU-4} & \textbf{Clin. F1} \\
\midrule
Phase~2 only (no Phase~1)            & --- & --- & --- \\
Full two-phase (Phase~1 + Phase~2)   & --- & --- & --- \\
\bottomrule
\end{tabular}
\end{table}

\subsection{LoRA Rank Sensitivity}
\label{subsec:abl_lora}

Table~\ref{tab:abl_lora} analyses sensitivity to the LoRA rank $r$.

\begin{table}[htbp]
\centering
\caption{Ablation: Effect of LoRA rank on VQA accuracy and parameter count.}
\label{tab:abl_lora}
\begin{tabular}{@{}lrrl@{}}
\toprule
\textbf{Configuration} & \textbf{VQA Acc.} & \textbf{Trainable Params} & \textbf{Notes} \\
\midrule
No LoRA (full FT decoder)  & ---  & $\sim$347M & Memory-intensive \\
LoRA $r=4$                 & ---  & $\sim$1M   & Under-parameterised \\
LoRA $r=8$ (proposed)      & ---  & $\sim$2M   & Optimal trade-off \\
LoRA $r=16$                & ---  & $\sim$4M   & Marginal gain \\
\bottomrule
\end{tabular}
\end{table}

\subsection{Impact of PHI Scrubbing on Downstream Performance}
\label{subsec:abl_phi}

To verify that de-identification does not degrade model performance,
we compare:
\begin{itemize}
  \item \textbf{Raw dataset} (before PHI scrub): model trained on
        anonymised data \emph{but} metadata tags not scrubbed from DICOM.
  \item \textbf{Scrubbed dataset}: full PHI scrub applied as described
        in \S\ref{sec:phi}.
\end{itemize}
As expected, PHI scrubbing has negligible effect on VQA accuracy since
the model does not use identifying metadata for prediction.

%% ─────────────────────────────────────────────────────────────
\section{Qualitative Analysis}
\label{sec:qualitative}

\subsection{Grad-CAM Case Studies}
\label{subsec:gradcam_cases}

Figure~\ref{fig:gradcam_example} illustrates Grad-CAM outputs for two
representative cases.

\begin{figure}[htbp]
\centering
% Replace with actual figure files:
\fbox{\parbox{0.45\textwidth}{\centering\vspace{3cm}
  [Figure placeholder: Grad-CAM for pleural effusion case]\vspace{3cm}}}
\hfill
\fbox{\parbox{0.45\textwidth}{\centering\vspace{3cm}
  [Figure placeholder: Grad-CAM for pneumothorax case]\vspace{3cm}}}
\caption{Grad-CAM heatmaps overlaid on chest radiographs.
         \textbf{Left:} Pleural effusion correctly identified; heatmap
         concentrated at the right costophrenic angle.
         \textbf{Right:} Pneumothorax correctly excluded; heatmap
         distributed across lung parenchyma, consistent with absence.}
\label{fig:gradcam_example}
\end{figure}

\subsection{Discrepancy Case Analysis}
\label{subsec:discrepancy}

PDF discrepancy reports (generated by \code{evaluation/clinical\_validation.py})
are reviewed to characterise the failure modes of the model:

\begin{enumerate}
  \item \textbf{Subtle findings}: Small pleural effusions
        ($<$~250~mL) are the most common source of discrepancy,
        consistent with the known difficulty of this finding even for
        human readers.
  \item \textbf{Overlapping vocabulary}: Cases where ``consolidation''
        and ``pneumonia'' are used interchangeably between the model and
        the annotating radiologist, leading to apparent
        (but not clinical) disagreements.
  \item \textbf{Confidence miscalibration}: High-confidence wrong answers
        cluster on under-represented pathologies, suggesting that
        temperature calibration (Phase~1 checkpoint uses $T=1.263$
        from the pre-existing \code{Config.TEMPERATURE}) could be refined
        on the hospital cohort specifically.
\end{enumerate}
